% ch6-produits-scalaires.tex — Chapitre 6 : Espaces préhilbertiens et orthogonalité
% Ce chapitre développe la théorie des produits scalaires en dimension finie :
% formes bilinéaires, formes quadratiques, produit scalaire, norme,
% orthogonalité, procédé de Gram-Schmidt, matrices orthogonales et unitaires,
% opérateur adjoint, et applications.

\chapter{Espaces préhilbertiens et orthogonalité}\label{ch:produits-scalaires}

\section*{Introduction}
\addcontentsline{toc}{section}{Introduction}

Les chapitres précédents ont construit la théorie des espaces vectoriels,
des applications linéaires, du déterminant et de la diagonalisation.
Cependant, cette théorie purement algébrique ne permet pas de définir des
notions fondamentales de la géométrie~: \emph{longueur} d'un vecteur,
\emph{angle} entre deux vecteurs, \emph{distance} entre deux points,
ou encore \emph{perpendicularité}\index{perpendicularité}.

\medskip

Pour introduire ces notions géométriques dans un espace vectoriel, il
faut le munir d'une structure supplémentaire~: un \emph{produit
scalaire}\index{produit scalaire}. C'est une forme bilinéaire (ou
sesquilinéaire dans le cas complexe) qui généralise le produit scalaire
usuel de $\R^n$. L'espace vectoriel muni d'un produit scalaire est
appelé \emph{espace préhilbertien}\index{espace préhilbertien}, ou
\emph{espace euclidien}\index{espace euclidien} lorsqu'il est réel et de
dimension finie, et \emph{espace hermitien}\index{espace hermitien}
lorsqu'il est complexe et de dimension finie.

\medskip

Ce chapitre est au c\oe ur de nombreuses applications~:
\begin{itemize}
  \item \textbf{Projection orthogonale~:}\index{projection orthogonale}
    trouver le vecteur d'un sous-espace le plus proche d'un vecteur donné
    (approximation au sens des moindres carrés).
  \item \textbf{Procédé de Gram-Schmidt~:}\index{Gram-Schmidt}
    construire une base orthonormée à partir d'une base quelconque, ce qui
    simplifie considérablement les calculs.
  \item \textbf{Analyse de Fourier~:}\index{Fourier}
    décomposer une fonction en somme de fonctions trigonométriques
    orthogonales.
  \item \textbf{Factorisation QR~:}\index{factorisation QR}
    décomposer une matrice en produit d'une matrice orthogonale et
    d'une matrice triangulaire, outil fondamental du calcul numérique.
  \item \textbf{Théorème spectral~:}\index{théorème spectral}
    les matrices symétriques réelles (ou hermitiennes) sont
    diagonalisables en base orthonormée, résultat essentiel pour
    l'analyse en composantes principales.
\end{itemize}

\medskip

Dans tout ce chapitre, $\K$ désigne $\R$ ou $\C$, $E$ un $\K$-espace
vectoriel, et $n=\dim E$ lorsque $E$ est de dimension finie.

% ============================================================
\section{Formes bilinéaires}\label{sec:formes-bilineaires}
% ============================================================

\begin{definition}{Forme bilinéaire}{def:forme-bilineaire}
  \index{forme bilinéaire}
  Soit $E$ un $\R$-espace vectoriel. Une \emph{forme bilinéaire} sur $E$
  est une application $\vphi\colon E\times E\to\R$ qui est linéaire en
  chacune de ses variables~:
  \begin{enumerate}[label=(\roman*)]
    \item Pour tout $\vec{v}\in E$, l'application
      $\vec{u}\mapsto\vphi(\vec{u},\vec{v})$ est linéaire.
    \item Pour tout $\vec{u}\in E$, l'application
      $\vec{v}\mapsto\vphi(\vec{u},\vec{v})$ est linéaire.
  \end{enumerate}
  Autrement dit, pour tous $\vec{u},\vec{u}',\vec{v}\in E$ et tout
  $\lambda\in\R$~:
  \[
    \vphi(\vec{u}+\vec{u}',\vec{v})=\vphi(\vec{u},\vec{v})
    +\vphi(\vec{u}',\vec{v}),\qquad
    \vphi(\lambda\vec{u},\vec{v})=\lambda\,\vphi(\vec{u},\vec{v}),
  \]
  et de même en la seconde variable.
\end{definition}

\begin{definition}{Symétrie, antisymétrie}{def:symetrie-bilineaire}
  \index{forme bilinéaire!symétrique}\index{forme bilinéaire!antisymétrique}
  Soit $\vphi$ une forme bilinéaire sur~$E$.
  \begin{enumerate}[label=(\roman*)]
    \item $\vphi$ est \emph{symétrique} si
      $\vphi(\vec{u},\vec{v})=\vphi(\vec{v},\vec{u})$ pour tous
      $\vec{u},\vec{v}\in E$.
    \item $\vphi$ est \emph{antisymétrique} si
      $\vphi(\vec{u},\vec{v})=-\vphi(\vec{v},\vec{u})$ pour tous
      $\vec{u},\vec{v}\in E$.
  \end{enumerate}
\end{definition}

\begin{definition}{Matrice d'une forme bilinéaire}{def:mat-bilineaire}
  \index{forme bilinéaire!matrice}
  Soit $\cB=(\vece{1},\ldots,\vece{n})$ une base de $E$ et $\vphi$ une
  forme bilinéaire sur $E$. La \emph{matrice de $\vphi$ dans la
  base~$\cB$} est la matrice $A=(a_{ij})\in\Mat_n(\R)$ définie par~:
  \[
    a_{ij}\deff\vphi(\vece{i},\vece{j}),\qquad 1\leq i,j\leq n.
  \]
  Pour tous $\vec{u},\vec{v}\in E$ de coordonnées $X,Y\in\R^n$ dans
  $\cB$, on a~:
  \[
    \vphi(\vec{u},\vec{v})=\transpose{X}\,A\,Y
    =\sum_{i=1}^n\sum_{j=1}^n a_{ij}\,x_i\,y_j.
  \]
\end{definition}

\begin{proposition}{Changement de base pour les formes bilinéaires}{prop:changement-base-bilineaire}
  \index{forme bilinéaire!changement de base}
  Si $P$ est la matrice de passage de la base $\cB$ à la base $\cB'$,
  et si $A$ (resp.\ $A'$) est la matrice de $\vphi$ dans $\cB$
  (resp.\ $\cB'$), alors~:
  \[
    A'=\transpose{P}\,A\,P.
  \]
  On dit que $A$ et $A'$ sont \emph{congruentes}\index{congruence de matrices}.
\end{proposition}

\begin{proof}
  Soient $\vec{u},\vec{v}\in E$ de coordonnées $X,Y$ dans $\cB$ et
  $X',Y'$ dans $\cB'$. On a $X=PX'$ et $Y=PY'$, donc~:
  \[
    \vphi(\vec{u},\vec{v})=\transpose{X}\,A\,Y
    =\transpose{(PX')}\,A\,(PY')
    =\transpose{X'}\,(\transpose{P}\,A\,P)\,Y'.
  \]
  Par unicité de la matrice représentant $\vphi$ dans $\cB'$, on obtient
  $A'=\transpose{P}\,A\,P$.
\end{proof}

\begin{remark}{Forme bilinéaire symétrique et matrice symétrique}{rem:sym-mat}
  $\vphi$ est symétrique si et seulement si $A=\transpose{A}$.
  En effet, $\vphi(\vec{u},\vec{v})=\transpose{X}AY$ et
  $\vphi(\vec{v},\vec{u})=\transpose{Y}AX=\transpose{X}\transpose{A}Y$.
  L'égalité pour tous $X,Y$ entraîne $A=\transpose{A}$.
\end{remark}

% ============================================================
\section{Formes quadratiques}\label{sec:formes-quadratiques}
% ============================================================

\begin{definition}{Forme quadratique}{def:forme-quadratique}
  \index{forme quadratique}
  Soit $E$ un $\R$-espace vectoriel. Une \emph{forme quadratique} sur $E$
  est une application $q\colon E\to\R$ telle qu'il existe une forme
  bilinéaire symétrique $\vphi$ sur $E$ vérifiant~:
  \[
    q(\vec{v})=\vphi(\vec{v},\vec{v})\quad
    \text{pour tout }\vec{v}\in E.
  \]
  On dit que $\vphi$ est la \emph{forme polaire}\index{forme polaire} de $q$.
\end{definition}

\begin{proposition}{Identité de polarisation}{prop:polarisation}
  \index{identité de polarisation}
  La forme bilinéaire symétrique $\vphi$ est entièrement déterminée
  par la forme quadratique $q$, via la \emph{formule de polarisation}~:
  \[
    \vphi(\vec{u},\vec{v})
    =\frac{1}{2}\bigl[q(\vec{u}+\vec{v})-q(\vec{u})-q(\vec{v})\bigr]
    =\frac{1}{4}\bigl[q(\vec{u}+\vec{v})-q(\vec{u}-\vec{v})\bigr].
  \]
\end{proposition}

\begin{proof}
  Par bilinéarité et symétrie de $\vphi$~:
  \begin{align*}
    q(\vec{u}+\vec{v})
    &=\vphi(\vec{u}+\vec{v},\vec{u}+\vec{v}) \\
    &=\vphi(\vec{u},\vec{u})+\vphi(\vec{u},\vec{v})
      +\vphi(\vec{v},\vec{u})+\vphi(\vec{v},\vec{v}) \\
    &=q(\vec{u})+2\,\vphi(\vec{u},\vec{v})+q(\vec{v}).
  \end{align*}
  D'où $\vphi(\vec{u},\vec{v})
  =\frac{1}{2}[q(\vec{u}+\vec{v})-q(\vec{u})-q(\vec{v})]$.
  De même, $q(\vec{u}-\vec{v})=q(\vec{u})-2\,\vphi(\vec{u},\vec{v})
  +q(\vec{v})$, ce qui donne la seconde formule par soustraction.
\end{proof}

\begin{theorem}{Loi d'inertie de Sylvester}{thm:sylvester}
  \index{loi d'inertie de Sylvester}\index{Sylvester}
  Soit $q$ une forme quadratique sur un $\R$-espace vectoriel $E$ de
  dimension $n$. Il existe une base $\cB$ de $E$ dans laquelle la
  matrice de la forme polaire $\vphi$ est diagonale~:
  \[
    \Mat_\cB(\vphi)=\diag(\underbrace{1,\ldots,1}_{p},\,
    \underbrace{-1,\ldots,-1}_{q},\,
    \underbrace{0,\ldots,0}_{n-p-q}).
  \]
  De plus, les entiers $p$ (indice de positivité) et $q$ (indice de
  négativité) ne dépendent pas du choix de la base~: ils sont des
  invariants de $q$. Le couple $(p,q)$ s'appelle la
  \emph{signature}\index{signature} de $q$, et $p+q$ est le
  \emph{rang}\index{rang!d'une forme quadratique} de $q$.
\end{theorem}

\begin{remark}{Interprétation de la signature}{rem:signature}
  La signature $(p,q)$ mesure le nombre de directions dans lesquelles $q$
  est strictement positive ($p$) ou strictement négative ($q$). Une forme
  quadratique est~:
  \begin{itemize}
    \item \emph{définie positive}\index{forme quadratique!définie positive}
      si $(p,q)=(n,0)$, i.e.\ $q(\vec{v})>0$ pour tout
      $\vec{v}\neq\vzero$.
    \item \emph{définie négative} si $(p,q)=(0,n)$.
    \item \emph{positive} (ou semi-définie positive) si $q=0$, i.e.\
      $q(\vec{v})\geq 0$ pour tout $\vec{v}$.
    \item \emph{non dégénérée}\index{forme quadratique!non dégénérée}
      si $p+q=n$.
  \end{itemize}
\end{remark}

% ============================================================
\section{Produit scalaire}\label{sec:produit-scalaire}
% ============================================================

\subsection{Cas réel}\label{subsec:ps-reel}

\begin{definition}{Produit scalaire réel}{def:ps-reel}
  \index{produit scalaire!réel}
  Soit $E$ un $\R$-espace vectoriel. Un \emph{produit scalaire} sur $E$
  est une forme bilinéaire $\inner{\cdot,\cdot}\colon E\times E\to\R$
  qui est~:
  \begin{enumerate}[label=(\roman*)]
    \item \textbf{symétrique~:} $\inner{\vec{u},\vec{v}}
      =\inner{\vec{v},\vec{u}}$ pour tous $\vec{u},\vec{v}\in E$.
    \item \textbf{définie positive~:} $\inner{\vec{v},\vec{v}}\geq 0$
      pour tout $\vec{v}\in E$, et $\inner{\vec{v},\vec{v}}=0$ si et
      seulement si $\vec{v}=\vzero$.
  \end{enumerate}
  Un $\R$-espace vectoriel muni d'un produit scalaire est un
  \emph{espace préhilbertien réel}. S'il est de dimension finie, on
  l'appelle \emph{espace euclidien}\index{espace euclidien}.
\end{definition}

\subsection{Cas complexe}\label{subsec:ps-complexe}

\begin{definition}{Forme sesquilinéaire}{def:sesquilineaire}
  \index{forme sesquilinéaire}
  Soit $E$ un $\C$-espace vectoriel. Une \emph{forme sesquilinéaire}
  sur $E$ est une application $\vphi\colon E\times E\to\C$ qui est~:
  \begin{enumerate}[label=(\roman*)]
    \item linéaire en la seconde variable~:
      $\vphi(\vec{u},\lambda\vec{v}+\vec{v}')
      =\lambda\,\vphi(\vec{u},\vec{v})+\vphi(\vec{u},\vec{v}')$.
    \item semi-linéaire (antilinéaire) en la première variable~:
      $\vphi(\lambda\vec{u}+\vec{u}',\vec{v})
      =\bar{\lambda}\,\vphi(\vec{u},\vec{v})+\vphi(\vec{u}',\vec{v})$.
  \end{enumerate}
  (Convention physique~: la conjugaison porte sur le premier argument.)
\end{definition}

\begin{definition}{Produit scalaire hermitien}{def:ps-hermitien}
  \index{produit scalaire!hermitien}\index{produit hermitien}
  Soit $E$ un $\C$-espace vectoriel. Un \emph{produit scalaire hermitien}
  (ou \emph{produit hermitien}) sur $E$ est une forme sesquilinéaire
  $\inner{\cdot,\cdot}\colon E\times E\to\C$ qui est~:
  \begin{enumerate}[label=(\roman*)]
    \item \textbf{hermitienne~:} $\inner{\vec{u},\vec{v}}
      =\overline{\inner{\vec{v},\vec{u}}}$ pour tous $\vec{u},\vec{v}\in E$.
    \item \textbf{définie positive~:} $\inner{\vec{v},\vec{v}}\in\R$,
      $\inner{\vec{v},\vec{v}}\geq 0$ pour tout $\vec{v}\in E$, et
      $\inner{\vec{v},\vec{v}}=0$ si et seulement si $\vec{v}=\vzero$.
  \end{enumerate}
  Un $\C$-espace vectoriel de dimension finie muni d'un produit hermitien
  est un \emph{espace hermitien}\index{espace hermitien}.
\end{definition}

\begin{remark}{Convention}{rem:convention-ps}
  La condition hermitienne entraîne
  $\inner{\vec{v},\vec{v}}=\overline{\inner{\vec{v},\vec{v}}}$, donc
  $\inner{\vec{v},\vec{v}}\in\R$. De plus, dans le cas réel, la
  condition hermitienne se réduit à la symétrie. On peut donc énoncer
  les résultats pour $\K=\R$ ou $\C$ de manière unifiée.
\end{remark}

% ============================================================
\section{Exemples fondamentaux}\label{sec:exemples-ps}
% ============================================================

\begin{example}{Produit scalaire canonique de $\R^n$}{ex:ps-Rn}
  \index{produit scalaire!canonique sur $\R^n$}
  Sur $E=\R^n$, le \emph{produit scalaire canonique} est~:
  \[
    \inner{\vec{x},\vec{y}}
    =\transpose{\vec{x}}\,\vec{y}
    =\sum_{i=1}^n x_i\,y_i,\qquad
    \vec{x}=(x_1,\ldots,x_n),\;\vec{y}=(y_1,\ldots,y_n).
  \]
  Il est bilinéaire, symétrique, et défini positif car
  $\inner{\vec{x},\vec{x}}=\sum_{i=1}^n x_i^2\geq 0$, avec égalité si
  et seulement si $\vec{x}=\vzero$.
\end{example}

\begin{example}{Produit hermitien canonique de $\C^n$}{ex:ps-Cn}
  \index{produit scalaire!canonique sur $\C^n$}
  Sur $E=\C^n$, le \emph{produit hermitien canonique} est~:
  \[
    \inner{\vec{z},\vec{w}}
    =\hermit{\vec{z}}\,\vec{w}
    =\sum_{i=1}^n \bar{z}_i\,w_i.
  \]
  On vérifie immédiatement la sesquilinéarité, la condition hermitienne,
  et la positivité~:
  $\inner{\vec{z},\vec{z}}=\sum_{i=1}^n\abs{z_i}^2\geq 0$.
\end{example}

\begin{example}{Produit scalaire $L^2$ sur $\mathcal{C}([a,b],\R)$}{ex:ps-L2}
  \index{produit scalaire!$L^2$}
  Sur l'espace $E=\mathcal{C}([a,b],\R)$ des fonctions continues de
  $[a,b]$ dans $\R$, on pose~:
  \[
    \inner{f,g}=\int_a^b f(t)\,g(t)\,dt.
  \]
  C'est un produit scalaire~: la bilinéarité et la symétrie sont
  immédiates. Pour le caractère défini positif~:
  $\inner{f,f}=\int_a^b f(t)^2\,dt\geq 0$, et si
  $\inner{f,f}=0$, alors $f(t)^2\equiv 0$ par continuité, donc
  $f=0$.
\end{example}

\begin{example}{Produit scalaire à poids}{ex:ps-poids}
  \index{produit scalaire!à poids}
  Soit $w\colon[a,b]\to\R_{>0}$ une fonction continue strictement
  positive (fonction poids). Alors~:
  \[
    \inner{f,g}_w=\int_a^b w(t)\,f(t)\,g(t)\,dt
  \]
  est un produit scalaire sur $\mathcal{C}([a,b],\R)$.
  Ce type de produit scalaire apparaît naturellement dans la théorie
  des polynômes orthogonaux (Legendre, Tchebychev, Hermite, Laguerre).
  \index{polynômes orthogonaux}
\end{example}

% ============================================================
\section{Norme, distance et inégalité de Cauchy-Schwarz}%
\label{sec:norme-cauchy-schwarz}
% ============================================================

\begin{definition}{Norme associée à un produit scalaire}{def:norme}
  \index{norme}
  Soit $(E,\inner{\cdot,\cdot})$ un espace préhilbertien. La
  \emph{norme} associée au produit scalaire est l'application
  $\norm{\cdot}\colon E\to\R_+$ définie par~:
  \[
    \norm{\vec{v}}\deff\sqrt{\inner{\vec{v},\vec{v}}}.
  \]
  La \emph{distance}\index{distance} entre deux vecteurs est
  $d(\vec{u},\vec{v})\deff\norm{\vec{u}-\vec{v}}$.
\end{definition}

\begin{theorem}{Inégalité de Cauchy-Schwarz}{thm:cauchy-schwarz}
  \index{inégalité de Cauchy-Schwarz}\index{Cauchy-Schwarz}
  Soit $(E,\inner{\cdot,\cdot})$ un espace préhilbertien (réel ou
  complexe). Pour tous $\vec{u},\vec{v}\in E$~:
  \[
    \boxed{\abs{\inner{\vec{u},\vec{v}}}
    \leq\norm{\vec{u}}\cdot\norm{\vec{v}}.}
  \]
  De plus, l'égalité a lieu si et seulement si $\vec{u}$ et $\vec{v}$
  sont liés (i.e.\ l'un est multiple scalaire de l'autre).
\end{theorem}

\begin{proof}
  Si $\vec{v}=\vzero$, les deux membres valent $0$ et l'inégalité est
  triviale. Supposons donc $\vec{v}\neq\vzero$.

  \emph{Cas réel.} Pour tout $t\in\R$, considérons~:
  \[
    p(t)\deff\norm{\vec{u}-t\vec{v}}^2
    =\inner{\vec{u}-t\vec{v},\vec{u}-t\vec{v}}
    =\norm{\vec{u}}^2-2t\,\inner{\vec{u},\vec{v}}+t^2\norm{\vec{v}}^2.
  \]
  C'est un polynôme du second degré en $t$ à coefficients réels, avec
  coefficient dominant $\norm{\vec{v}}^2>0$, et $p(t)\geq 0$ pour tout
  $t\in\R$ (car c'est le carré d'une norme). La condition de
  non-négativité d'un trinôme du second degré impose que le discriminant
  soit négatif ou nul~:
  \[
    \Delta=4\,\inner{\vec{u},\vec{v}}^2
    -4\,\norm{\vec{u}}^2\norm{\vec{v}}^2\leq 0,
  \]
  d'où $\inner{\vec{u},\vec{v}}^2\leq\norm{\vec{u}}^2\norm{\vec{v}}^2$,
  ce qui est exactement l'inégalité de Cauchy-Schwarz.

  \emph{Cas d'égalité.} L'égalité a lieu si et seulement si $\Delta=0$,
  c'est-à-dire si et seulement si $p$ admet une racine réelle $t_0$,
  auquel cas $p(t_0)=\norm{\vec{u}-t_0\vec{v}}^2=0$, donc
  $\vec{u}=t_0\vec{v}$~: les vecteurs sont colinéaires.

  \emph{Cas complexe.} Pour tout $\lambda\in\C$, on a~:
  \[
    0\leq\norm{\vec{u}-\lambda\vec{v}}^2
    =\norm{\vec{u}}^2-\lambda\,\inner{\vec{u},\vec{v}}
    -\bar{\lambda}\,\overline{\inner{\vec{u},\vec{v}}}
    +\abs{\lambda}^2\norm{\vec{v}}^2.
  \]
  Choisissons $\lambda=\frac{\inner{\vec{u},\vec{v}}}{\norm{\vec{v}}^2}$
  (qui minimise l'expression). On obtient~:
  \[
    0\leq\norm{\vec{u}}^2
    -\frac{\abs{\inner{\vec{u},\vec{v}}}^2}{\norm{\vec{v}}^2},
  \]
  d'où $\abs{\inner{\vec{u},\vec{v}}}^2
  \leq\norm{\vec{u}}^2\cdot\norm{\vec{v}}^2$, et en prenant la racine
  carrée~:
  \[
    \abs{\inner{\vec{u},\vec{v}}}\leq\norm{\vec{u}}\cdot\norm{\vec{v}}.
  \]
  L'égalité a lieu si et seulement si
  $\norm{\vec{u}-\lambda\vec{v}}=0$, i.e.\
  $\vec{u}=\lambda\vec{v}$.
\end{proof}

\begin{corollary}{Inégalité triangulaire}{cor:ineg-triangulaire}
  \index{inégalité triangulaire}
  Pour tous $\vec{u},\vec{v}\in E$~:
  \[
    \norm{\vec{u}+\vec{v}}\leq\norm{\vec{u}}+\norm{\vec{v}}.
  \]
\end{corollary}

\begin{proof}
  On a~:
  \begin{align*}
    \norm{\vec{u}+\vec{v}}^2
    &=\norm{\vec{u}}^2+2\,\Re\bigl(\inner{\vec{u},\vec{v}}\bigr)
      +\norm{\vec{v}}^2 \\
    &\leq\norm{\vec{u}}^2+2\,\abs{\inner{\vec{u},\vec{v}}}
      +\norm{\vec{v}}^2 \\
    &\leq\norm{\vec{u}}^2+2\,\norm{\vec{u}}\norm{\vec{v}}
      +\norm{\vec{v}}^2
    =\bigl(\norm{\vec{u}}+\norm{\vec{v}}\bigr)^2,
  \end{align*}
  en utilisant Cauchy-Schwarz à la dernière inégalité. On conclut en
  prenant la racine carrée.
\end{proof}

\begin{proposition}{Propriétés de la norme}{prop:norme}
  La norme $\norm{\cdot}$ vérifie les propriétés suivantes~:
  \begin{enumerate}[label=(\roman*)]
    \item \textbf{Séparation~:} $\norm{\vec{v}}=0\Leftrightarrow
      \vec{v}=\vzero$.
    \item \textbf{Homogénéité~:}
      $\norm{\lambda\vec{v}}=\abs{\lambda}\,\norm{\vec{v}}$ pour tout
      $\lambda\in\K$.
    \item \textbf{Inégalité triangulaire~:}
      $\norm{\vec{u}+\vec{v}}\leq\norm{\vec{u}}+\norm{\vec{v}}$.
  \end{enumerate}
  De plus, elle satisfait l'\emph{identité du parallélogramme}%
  \index{identité du parallélogramme}~:
  \[
    \norm{\vec{u}+\vec{v}}^2+\norm{\vec{u}-\vec{v}}^2
    =2\bigl(\norm{\vec{u}}^2+\norm{\vec{v}}^2\bigr).
  \]
\end{proposition}

\begin{definition}{Angle entre deux vecteurs}{def:angle}
  \index{angle}
  Soient $\vec{u},\vec{v}\in E\setminus\set{\vzero}$ dans un espace
  euclidien. L'inégalité de Cauchy-Schwarz assure que~:
  \[
    -1\leq\frac{\inner{\vec{u},\vec{v}}}%
    {\norm{\vec{u}}\cdot\norm{\vec{v}}}\leq 1.
  \]
  On définit l'\emph{angle} $\theta\in[0,\pi]$ entre $\vec{u}$ et
  $\vec{v}$ par~:
  \[
    \cos\theta=\frac{\inner{\vec{u},\vec{v}}}%
    {\norm{\vec{u}}\cdot\norm{\vec{v}}}.
  \]
\end{definition}

% ============================================================
\section{Orthogonalité}\label{sec:orthogonalite}
% ============================================================

\begin{definition}{Vecteurs orthogonaux}{def:ortho}
  \index{orthogonalité}\index{vecteurs orthogonaux}
  Soit $(E,\inner{\cdot,\cdot})$ un espace préhilbertien.
  \begin{enumerate}[label=(\roman*)]
    \item Deux vecteurs $\vec{u},\vec{v}\in E$ sont \emph{orthogonaux},
      noté $\vec{u}\perp\vec{v}$, si $\inner{\vec{u},\vec{v}}=0$.
    \item Une famille $(\vec{v}_1,\ldots,\vec{v}_p)$ est
      \emph{orthogonale}\index{famille orthogonale} si
      $\inner{\vec{v}_i,\vec{v}_j}=0$ pour tout $i\neq j$.
    \item Elle est \emph{orthonormée}\index{famille orthonormée}
      (ou \emph{orthonormale}) si de plus
      $\norm{\vec{v}_i}=1$ pour tout $i$, i.e.\
      $\inner{\vec{v}_i,\vec{v}_j}=\delta_{ij}$.
  \end{enumerate}
\end{definition}

\begin{proposition}{Famille orthogonale et liberté}{prop:ortho-libre}
  \index{famille orthogonale!liberté}
  Toute famille orthogonale de vecteurs \emph{non nuls} est libre.
\end{proposition}

\begin{proof}
  Soit $(\vec{v}_1,\ldots,\vec{v}_p)$ une famille orthogonale avec
  $\vec{v}_i\neq\vzero$ pour tout $i$. Supposons que
  $\sum_{i=1}^p\alpha_i\vec{v}_i=\vzero$. En prenant le produit
  scalaire avec $\vec{v}_k$~:
  \[
    0=\inner*{\sum_{i=1}^p\alpha_i\vec{v}_i,\vec{v}_k}
    =\sum_{i=1}^p\alpha_i\inner{\vec{v}_i,\vec{v}_k}
    =\alpha_k\norm{\vec{v}_k}^2.
  \]
  Comme $\norm{\vec{v}_k}^2>0$ (car $\vec{v}_k\neq\vzero$), on en
  déduit $\alpha_k=0$ pour tout $k$. La famille est donc libre.
\end{proof}

\begin{definition}{Orthogonal d'une partie}{def:orthogonal}
  \index{orthogonal d'un sous-espace}\index{$F^\perp$}
  Soit $F$ une partie de $E$. L'\emph{orthogonal} de $F$ est~:
  \[
    F^\perp\deff\setcond{\vec{v}\in E}%
    {\inner{\vec{v},\vec{w}}=0\;\text{pour tout }\vec{w}\in F}.
  \]
\end{definition}

\begin{proposition}{Propriétés de l'orthogonal}{prop:orthogonal}
  Soient $F,G$ des sous-espaces vectoriels d'un espace préhilbertien
  $E$ de dimension finie. Alors~:
  \begin{enumerate}[label=(\roman*)]
    \item $F^\perp$ est un sous-espace vectoriel de $E$.
    \item $E=F\oplus F^\perp$.
      \label{item:decomp-ortho}
    \item $\dim F^\perp=\dim E-\dim F$.
    \item $(F^\perp)^\perp=F$.
    \item $F\subset G\implies G^\perp\subset F^\perp$.
    \item $(F+G)^\perp=F^\perp\cap G^\perp$.
    \item $(F\cap G)^\perp=F^\perp+G^\perp$.
  \end{enumerate}
\end{proposition}

\begin{proof}
  \emph{(i)} La linéarité du produit scalaire en chaque variable assure
  que $F^\perp$ est stable par combinaison linéaire.

  \emph{(ii)} Montrons d'abord que $F\cap F^\perp=\set{\vzero}$~:
  si $\vec{v}\in F\cap F^\perp$, alors
  $\inner{\vec{v},\vec{v}}=0$ (car $\vec{v}\in F^\perp$ et
  $\vec{v}\in F$), donc $\vec{v}=\vzero$.

  Soit $(\vece{1},\ldots,\vece{p})$ une base orthonormée de $F$
  (dont l'existence sera établie par le procédé de Gram-Schmidt,
  
ef{thm:thm:gram-schmidt}). Pour tout $\vec{v}\in E$, posons~:
  \[
    \vec{v}_F\deff\sum_{i=1}^p\inner{\vece{i},\vec{v}}\,\vece{i},
    \qquad
    \vec{v}_{F^\perp}\deff\vec{v}-\vec{v}_F.
  \]
  On a $\vec{v}_F\in F$. Vérifions que $\vec{v}_{F^\perp}\in F^\perp$~:
  pour tout $j\in\set{1,\ldots,p}$,
  \[
    \inner{\vece{j},\vec{v}_{F^\perp}}
    =\inner{\vece{j},\vec{v}}
    -\sum_{i=1}^p\inner{\vece{i},\vec{v}}\,\inner{\vece{j},\vece{i}}
    =\inner{\vece{j},\vec{v}}-\inner{\vece{j},\vec{v}}=0.
  \]
  Donc $\vec{v}=\vec{v}_F+\vec{v}_{F^\perp}\in F+F^\perp$. Avec
  $F\cap F^\perp=\set{\vzero}$, on obtient $E=F\oplus F^\perp$.

  \emph{(iii)} Conséquence immédiate de (ii)~:
  $\dim F^\perp=\dim E-\dim F$.

  \emph{(iv)} L'inclusion $F\subset(F^\perp)^\perp$ est claire (si
  $\vec{v}\in F$, alors $\inner{\vec{v},\vec{w}}=0$ pour tout
  $\vec{w}\in F^\perp$). L'égalité des dimensions
  $\dim(F^\perp)^\perp=n-\dim F^\perp=n-(n-\dim F)=\dim F$ donne
  l'égalité.

  Les propriétés (v)--(vii) se vérifient par double inclusion à partir
  des définitions.
\end{proof}

\begin{theorem}{Théorème de Pythagore}{thm:pythagore}
  \index{théorème de Pythagore}
  Soient $\vec{u},\vec{v}\in E$ deux vecteurs orthogonaux. Alors~:
  \[
    \norm{\vec{u}+\vec{v}}^2=\norm{\vec{u}}^2+\norm{\vec{v}}^2.
  \]
  Plus généralement, si $(\vec{v}_1,\ldots,\vec{v}_p)$ est une famille
  orthogonale, alors~:
  \[
    \norm*{\sum_{i=1}^p\vec{v}_i}^2=\sum_{i=1}^p\norm{\vec{v}_i}^2.
  \]
\end{theorem}

\begin{proof}
  $\norm{\vec{u}+\vec{v}}^2
  =\inner{\vec{u}+\vec{v},\vec{u}+\vec{v}}
  =\norm{\vec{u}}^2+\inner{\vec{u},\vec{v}}
  +\inner{\vec{v},\vec{u}}+\norm{\vec{v}}^2
  =\norm{\vec{u}}^2+\norm{\vec{v}}^2$,
  puisque $\inner{\vec{u},\vec{v}}=0$ entraîne aussi
  $\inner{\vec{v},\vec{u}}=\overline{\inner{\vec{u},\vec{v}}}=0$.
  La généralisation se fait par récurrence immédiate.
\end{proof}

% ============================================================
\section{Projection orthogonale}\label{sec:projection-orthogonale}
% ============================================================

\begin{definition}{Projection orthogonale}{def:proj-ortho}
  \index{projection orthogonale}
  Soit $F$ un sous-espace vectoriel d'un espace préhilbertien $E$ de
  dimension finie. D'après la décomposition $E=F\oplus F^\perp$, tout
  vecteur $\vec{v}\in E$ s'écrit de manière unique~:
  \[
    \vec{v}=\vec{v}_F+\vec{v}_{F^\perp},\qquad
    \vec{v}_F\in F,\quad\vec{v}_{F^\perp}\in F^\perp.
  \]
  L'application $\proj_F\colon E\to E$ définie par
  $\proj_F(\vec{v})\deff\vec{v}_F$ est la \emph{projection orthogonale}
  sur $F$ parallèlement à $F^\perp$.
\end{definition}

\begin{theorem}{Caractérisation de la meilleure approximation}%
{thm:meilleure-approx}
  \index{meilleure approximation}
  Soit $F$ un sous-espace vectoriel de dimension finie d'un espace
  préhilbertien $E$, et soit $\vec{v}\in E$. L'élément
  $\vec{v}_F=\proj_F(\vec{v})$ est l'unique élément de $F$ qui
  minimise la distance de $\vec{v}$ à $F$~:
  \[
    \norm{\vec{v}-\vec{v}_F}
    =\min_{\vec{w}\in F}\norm{\vec{v}-\vec{w}}
    =d(\vec{v},F).
  \]
  De plus, $\vec{v}_F$ est caractérisé par les deux conditions~:
  \begin{enumerate}[label=(\roman*)]
    \item $\vec{v}_F\in F$.
    \item $\vec{v}-\vec{v}_F\perp F$, i.e.\
      $\inner{\vec{v}-\vec{v}_F,\vec{w}}=0$ pour tout $\vec{w}\in F$.
  \end{enumerate}
\end{theorem}

\begin{proof}
  Pour tout $\vec{w}\in F$, on a
  $\vec{v}-\vec{w}=(\vec{v}-\vec{v}_F)+(\vec{v}_F-\vec{w})$, avec
  $\vec{v}-\vec{v}_F\in F^\perp$ et
  $\vec{v}_F-\vec{w}\in F$. Ces deux vecteurs sont orthogonaux, donc
  par le théorème de Pythagore~:
  \[
    \norm{\vec{v}-\vec{w}}^2
    =\norm{\vec{v}-\vec{v}_F}^2+\norm{\vec{v}_F-\vec{w}}^2
    \geq\norm{\vec{v}-\vec{v}_F}^2.
  \]
  L'égalité a lieu si et seulement si
  $\norm{\vec{v}_F-\vec{w}}=0$, i.e.\ $\vec{w}=\vec{v}_F$.
  L'unicité est ainsi prouvée.
\end{proof}

\begin{proposition}{Formule de la projection orthogonale}{prop:formule-proj}
  \index{projection orthogonale!formule}
  Si $(\vece{1},\ldots,\vece{p})$ est une base \emph{orthonormée}
  de $F$, alors~:
  \[
    \proj_F(\vec{v})
    =\sum_{i=1}^p\inner{\vece{i},\vec{v}}\,\vece{i}.
  \]
\end{proposition}

\begin{proof}
  Posons $\vec{v}_F=\sum_{i=1}^p\inner{\vece{i},\vec{v}}\,\vece{i}
  \in F$. Pour tout $j$~:
  \[
    \inner{\vece{j},\vec{v}-\vec{v}_F}
    =\inner{\vece{j},\vec{v}}
    -\sum_{i=1}^p\inner{\vece{i},\vec{v}}\,\delta_{ij}
    =0.
  \]
  Donc $\vec{v}-\vec{v}_F\in F^\perp$, et par unicité de la
  décomposition $E=F\oplus F^\perp$, $\vec{v}_F=\proj_F(\vec{v})$.
\end{proof}

% ============================================================
\section{Procédé de Gram-Schmidt}\label{sec:gram-schmidt}
% ============================================================

\begin{theorem}{Procédé d'orthonormalisation de Gram-Schmidt}%
{thm:gram-schmidt}
  \index{Gram-Schmidt}\index{procédé de Gram-Schmidt}
  Soit $(\vec{v}_1,\ldots,\vec{v}_p)$ une famille libre dans un espace
  préhilbertien $E$. Il existe une famille orthonormée
  $(\vece{1},\ldots,\vece{p})$ telle que~:
  \[
    \Vect(\vece{1},\ldots,\vece{k})
    =\Vect(\vec{v}_1,\ldots,\vec{v}_k)
    \quad\text{pour tout }k\in\set{1,\ldots,p}.
  \]
  De plus, on peut choisir les $\vece{k}$ de sorte que
  $\inner{\vece{k},\vec{v}_k}>0$.
\end{theorem}

\begin{proof}
  On procède par récurrence sur~$k$.

  \emph{Initialisation ($k=1$).} On pose
  $\vece{1}\deff\frac{\vec{v}_1}{\norm{\vec{v}_1}}$
  (bien défini car $\vec{v}_1\neq\vzero$ puisque la famille est libre).
  On a $\norm{\vece{1}}=1$ et
  $\Vect(\vece{1})=\Vect(\vec{v}_1)$.

  \emph{Hérédité.} Supposons construits
  $\vece{1},\ldots,\vece{k-1}$ orthonormés avec
  $\Vect(\vece{1},\ldots,\vece{j})
  =\Vect(\vec{v}_1,\ldots,\vec{v}_j)$ pour $j\leq k-1$.
  Posons $F_{k-1}=\Vect(\vece{1},\ldots,\vece{k-1})$ et~:
  \[
    \vec{w}_k\deff\vec{v}_k
    -\sum_{i=1}^{k-1}\inner{\vece{i},\vec{v}_k}\,\vece{i}
    =\vec{v}_k-\proj_{F_{k-1}}(\vec{v}_k).
  \]
  Le vecteur $\vec{w}_k$ est la composante de $\vec{v}_k$ orthogonale
  à $F_{k-1}$. Comme $\vec{v}_k\notin F_{k-1}=\Vect(\vec{v}_1,
  \ldots,\vec{v}_{k-1})$ (la famille $(\vec{v}_i)$ étant libre),
  on a $\vec{w}_k\neq\vzero$.

  Vérifions l'orthogonalité~: pour tout $j\in\set{1,\ldots,k-1}$,
  \[
    \inner{\vece{j},\vec{w}_k}
    =\inner{\vece{j},\vec{v}_k}
    -\sum_{i=1}^{k-1}\inner{\vece{i},\vec{v}_k}\,
    \underbrace{\inner{\vece{j},\vece{i}}}_{=\delta_{ij}}
    =\inner{\vece{j},\vec{v}_k}-\inner{\vece{j},\vec{v}_k}=0.
  \]

  On pose alors $\vece{k}\deff\frac{\vec{w}_k}{\norm{\vec{w}_k}}$,
  qui vérifie $\norm{\vece{k}}=1$ et $\vece{k}\perp\vece{j}$ pour
  $j<k$. De plus~:
  \[
    \vec{v}_k=\vec{w}_k+\sum_{i=1}^{k-1}\inner{\vece{i},\vec{v}_k}
    \,\vece{i}
    \in\Vect(\vece{1},\ldots,\vece{k}),
  \]
  donc $\Vect(\vec{v}_1,\ldots,\vec{v}_k)
  \subset\Vect(\vece{1},\ldots,\vece{k})$. Comme $\vece{k}$ est
  combinaison linéaire de $\vec{v}_k$ et des $\vece{i}$ ($i<k$),
  et que $\Vect(\vece{1},\ldots,\vece{k-1})
  =\Vect(\vec{v}_1,\ldots,\vec{v}_{k-1})$ par hypothèse de
  récurrence, l'inclusion réciproque est vérifiée. D'où l'égalité des
  sous-espaces.
\end{proof}

\begin{example}{Gram-Schmidt dans $\R^3$}{ex:gram-schmidt-R3}
  \index{Gram-Schmidt!exemple}
  Orthonormalisons la base $(\vec{v}_1,\vec{v}_2,\vec{v}_3)$ de $\R^3$
  avec le produit scalaire canonique, où~:
  \[
    \vec{v}_1=\begin{pmatrix}1\\1\\0\end{pmatrix},\quad
    \vec{v}_2=\begin{pmatrix}1\\0\\1\end{pmatrix},\quad
    \vec{v}_3=\begin{pmatrix}0\\1\\1\end{pmatrix}.
  \]

  \emph{Étape 1.}
  $\norm{\vec{v}_1}=\sqrt{1+1+0}=\sqrt{2}$, donc
  $\vece{1}=\frac{1}{\sqrt{2}}\begin{pmatrix}1\\1\\0\end{pmatrix}$.

  \emph{Étape 2.}
  $\inner{\vece{1},\vec{v}_2}
  =\frac{1}{\sqrt{2}}(1\cdot 1+1\cdot 0+0\cdot 1)=\frac{1}{\sqrt{2}}$.

  $\vec{w}_2=\vec{v}_2-\inner{\vece{1},\vec{v}_2}\,\vece{1}
  =\begin{pmatrix}1\\0\\1\end{pmatrix}
  -\frac{1}{\sqrt{2}}\cdot\frac{1}{\sqrt{2}}
  \begin{pmatrix}1\\1\\0\end{pmatrix}
  =\begin{pmatrix}1\\0\\1\end{pmatrix}
  -\begin{pmatrix}1/2\\1/2\\0\end{pmatrix}
  =\begin{pmatrix}1/2\\-1/2\\1\end{pmatrix}$.

  $\norm{\vec{w}_2}=\sqrt{1/4+1/4+1}=\sqrt{3/2}$, donc
  $\vece{2}=\sqrt{\frac{2}{3}}\begin{pmatrix}1/2\\-1/2\\1\end{pmatrix}
  =\frac{1}{\sqrt{6}}\begin{pmatrix}1\\-1\\2\end{pmatrix}$.

  \emph{Étape 3.}
  $\inner{\vece{1},\vec{v}_3}
  =\frac{1}{\sqrt{2}}(0+1+0)=\frac{1}{\sqrt{2}}$,\quad
  $\inner{\vece{2},\vec{v}_3}
  =\frac{1}{\sqrt{6}}(0-1+2)=\frac{1}{\sqrt{6}}$.

  $\vec{w}_3=\vec{v}_3
  -\inner{\vece{1},\vec{v}_3}\,\vece{1}
  -\inner{\vece{2},\vec{v}_3}\,\vece{2}
  =\begin{pmatrix}0\\1\\1\end{pmatrix}
  -\frac{1}{2}\begin{pmatrix}1\\1\\0\end{pmatrix}
  -\frac{1}{6}\begin{pmatrix}1\\-1\\2\end{pmatrix}
  =\begin{pmatrix}-2/3\\2/3\\2/3\end{pmatrix}$.

  $\norm{\vec{w}_3}=\frac{2}{3}\sqrt{3}$, donc
  $\vece{3}=\frac{1}{\sqrt{3}}\begin{pmatrix}-1\\1\\1\end{pmatrix}$.

  On peut vérifier que
  $\inner{\vece{i},\vece{j}}=\delta_{ij}$.
\end{example}

% ============================================================
\section{Bases orthonormées}\label{sec:bases-orthonormees}
% ============================================================

\begin{corollary}{Existence de bases orthonormées}{cor:existence-bon}
  \index{base orthonormée!existence}
  Tout espace préhilbertien de dimension finie admet une base
  orthonormée.
\end{corollary}

\begin{proof}
  Appliquer le procédé de Gram-Schmidt à une base quelconque.
\end{proof}

\begin{proposition}{Coordonnées dans une base orthonormée}{prop:coord-bon}
  \index{base orthonormée!coordonnées}
  Soit $(\vece{1},\ldots,\vece{n})$ une base orthonormée de $E$. Pour
  tout $\vec{v}\in E$~:
  \[
    \vec{v}=\sum_{i=1}^n\inner{\vece{i},\vec{v}}\,\vece{i}.
  \]
  Les scalaires $\inner{\vece{i},\vec{v}}$ sont les \emph{coefficients
  de Fourier}\index{coefficients de Fourier} de $\vec{v}$ dans la base
  $(\vece{i})$.
\end{proposition}

\begin{proof}
  En écrivant $\vec{v}=\sum_{i=1}^n\alpha_i\vece{i}$ et en prenant le
  produit scalaire avec $\vece{j}$, on obtient
  $\inner{\vece{j},\vec{v}}=\alpha_j$.
\end{proof}

\begin{theorem}{Identité de Parseval}{thm:parseval}
  \index{identité de Parseval}\index{Parseval}
  Soit $(\vece{1},\ldots,\vece{n})$ une base orthonormée de $E$.
  Pour tous $\vec{u},\vec{v}\in E$~:
  \[
    \inner{\vec{u},\vec{v}}
    =\sum_{i=1}^n\overline{\inner{\vece{i},\vec{u}}}\,
    \inner{\vece{i},\vec{v}}.
  \]
  En particulier~:
  \[
    \norm{\vec{v}}^2=\sum_{i=1}^n\abs{\inner{\vece{i},\vec{v}}}^2.
  \]
\end{theorem}

\begin{proof}
  En décomposant $\vec{u}=\sum_i\alpha_i\vece{i}$ et
  $\vec{v}=\sum_j\beta_j\vece{j}$ avec
  $\alpha_i=\inner{\vece{i},\vec{u}}$ et
  $\beta_j=\inner{\vece{j},\vec{v}}$~:
  \[
    \inner{\vec{u},\vec{v}}
    =\inner*{\sum_i\alpha_i\vece{i},\sum_j\beta_j\vece{j}}
    =\sum_i\sum_j\bar{\alpha}_i\beta_j
    \underbrace{\inner{\vece{i},\vece{j}}}_{=\delta_{ij}}
    =\sum_i\bar{\alpha}_i\beta_i.
    \qedhere
  \]
\end{proof}

\begin{theorem}{Inégalité de Bessel}{thm:bessel}
  \index{inégalité de Bessel}\index{Bessel}
  Soit $(\vece{1},\ldots,\vece{p})$ une famille orthonormée (non
  nécessairement base) de $E$. Pour tout $\vec{v}\in E$~:
  \[
    \sum_{i=1}^p\abs{\inner{\vece{i},\vec{v}}}^2
    \leq\norm{\vec{v}}^2.
  \]
  L'égalité a lieu si et seulement si $\vec{v}\in\Vect(\vece{1},
  \ldots,\vece{p})$.
\end{theorem}

\begin{proof}
  Posons $F=\Vect(\vece{1},\ldots,\vece{p})$ et
  $\vec{v}_F=\proj_F(\vec{v})=\sum_{i=1}^p\inner{\vece{i},\vec{v}}
  \,\vece{i}$. Par le théorème de Pythagore~:
  \[
    \norm{\vec{v}}^2=\norm{\vec{v}_F}^2+\norm{\vec{v}-\vec{v}_F}^2
    =\sum_{i=1}^p\abs{\inner{\vece{i},\vec{v}}}^2
    +\norm{\vec{v}-\vec{v}_F}^2
    \geq\sum_{i=1}^p\abs{\inner{\vece{i},\vec{v}}}^2.
  \]
  L'égalité a lieu ssi $\norm{\vec{v}-\vec{v}_F}=0$, i.e.\
  $\vec{v}=\vec{v}_F\in F$.
\end{proof}

% ============================================================
\section{Matrices orthogonales et unitaires}\label{sec:ortho-unitaire}
% ============================================================

\begin{definition}{Matrice orthogonale}{def:mat-ortho}
  \index{matrice orthogonale}\index{$\Orth(n)$}
  Une matrice $Q\in\Mat_n(\R)$ est \emph{orthogonale} si~:
  \[
    \transpose{Q}\,Q=I_n,
  \]
  ou de manière équivalente, $\inv{Q}=\transpose{Q}$. L'ensemble des
  matrices orthogonales de taille $n$ forme un groupe pour la
  multiplication, noté $\Orth(n)$ et appelé \emph{groupe orthogonal}%
  \index{groupe orthogonal}.
\end{definition}

\begin{definition}{Matrice unitaire}{def:mat-unitaire}
  \index{matrice unitaire}\index{$\mathrm{U}(n)$}
  Une matrice $U\in\Mat_n(\C)$ est \emph{unitaire} si~:
  \[
    \hermit{U}\,U=I_n,
  \]
  ou de manière équivalente, $\inv{U}=\hermit{U}$. L'ensemble des
  matrices unitaires de taille $n$ forme le \emph{groupe unitaire}
  $\mathrm{U}(n)$.
\end{definition}

\begin{proposition}{Caractérisations des matrices orthogonales}%
{prop:caract-ortho}
  Pour $Q\in\Mat_n(\R)$, les conditions suivantes sont équivalentes~:
  \begin{enumerate}[label=(\roman*)]
    \item $Q$ est orthogonale.
    \item Les colonnes de $Q$ forment une base orthonormée de $\R^n$.
    \item Les lignes de $Q$ forment une base orthonormée de $\R^n$.
    \item $Q$ conserve le produit scalaire~:
      $\inner{Q\vec{x},Q\vec{y}}=\inner{\vec{x},\vec{y}}$ pour tous
      $\vec{x},\vec{y}\in\R^n$.
    \item $Q$ est une isométrie~:
      $\norm{Q\vec{x}}=\norm{\vec{x}}$ pour tout $\vec{x}\in\R^n$.
  \end{enumerate}
\end{proposition}

\begin{proof}
  (i)$\Leftrightarrow$(ii)~: $\transpose{Q}Q=I_n$ signifie que le
  produit scalaire de la $i$-ème et de la $j$-ème colonne de $Q$ vaut
  $\delta_{ij}$.

  (ii)$\Leftrightarrow$(iii)~: Si $\transpose{Q}Q=I_n$, alors $Q$ est
  inversible et $Q\transpose{Q}=I_n$, ce qui signifie que les lignes
  sont orthonormées.

  (i)$\Rightarrow$(iv)~:
  $\inner{Q\vec{x},Q\vec{y}}=\transpose{(Q\vec{x})}Q\vec{y}
  =\transpose{\vec{x}}\transpose{Q}Q\vec{y}=\transpose{\vec{x}}\vec{y}
  =\inner{\vec{x},\vec{y}}$.

  (iv)$\Rightarrow$(v)~: prendre $\vec{y}=\vec{x}$.

  (v)$\Rightarrow$(i)~: par l'identité de polarisation,
  $\norm{Q\vec{x}}=\norm{\vec{x}}$ pour tout $\vec{x}$ entraîne
  $\inner{Q\vec{x},Q\vec{y}}=\inner{\vec{x},\vec{y}}$, d'où
  $\transpose{\vec{x}}(\transpose{Q}Q-I_n)\vec{y}=0$ pour tous
  $\vec{x},\vec{y}$, ce qui impose $\transpose{Q}Q=I_n$.
\end{proof}

\begin{proposition}{Déterminant d'une matrice orthogonale}{prop:det-ortho}
  \index{matrice orthogonale!déterminant}
  Si $Q\in\Orth(n)$, alors $\det(Q)=\pm 1$.
\end{proposition}

\begin{proof}
  $\det(\transpose{Q}Q)=\det(\transpose{Q})\det(Q)=\det(Q)^2=\det(I_n)=1$.
\end{proof}

\begin{definition}{Groupe spécial orthogonal}{def:SO}
  \index{$\SO(n)$}\index{groupe spécial orthogonal}
  Le \emph{groupe spécial orthogonal} est~:
  \[
    \SO(n)\deff\setcond{Q\in\Orth(n)}{\det(Q)=1}.
  \]
  Les éléments de $\SO(n)$ sont les \emph{rotations}%
  \index{rotation} de $\R^n$.
\end{definition}

\begin{example}{Matrices de $\SO(2)$}{ex:SO2}
  \index{$\SO(2)$}
  Les éléments de $\SO(2)$ sont les matrices de rotation d'angle
  $\theta$~:
  \[
    R_\theta=\begin{pmatrix}\cos\theta&-\sin\theta\\
    \sin\theta&\cos\theta\end{pmatrix},\qquad\theta\in\R.
  \]
  Les éléments de $\Orth(2)\setminus\SO(2)$ (de déterminant $-1$) sont
  les réflexions~:
  $S_\theta=\begin{pmatrix}\cos\theta&\sin\theta\\
  \sin\theta&-\cos\theta\end{pmatrix}$.
\end{example}

% ============================================================
\section{Opérateur adjoint}\label{sec:adjoint}
% ============================================================

\begin{definition}{Adjoint d'un endomorphisme}{def:adjoint}
  \index{adjoint}\index{opérateur adjoint}
  Soit $(E,\inner{\cdot,\cdot})$ un espace préhilbertien de dimension
  finie et $f\in\End(E)$. L'\emph{adjoint} de $f$ est l'unique
  endomorphisme $f^*\in\End(E)$ tel que~:
  \[
    \inner{f(\vec{u}),\vec{v}}=\inner{\vec{u},f^*(\vec{v})}
    \quad\text{pour tous }\vec{u},\vec{v}\in E.
  \]
\end{definition}

\begin{theorem}{Existence et unicité de l'adjoint}{thm:adjoint-existence}
  \index{adjoint!existence}
  Pour tout $f\in\End(E)$, l'adjoint $f^*$ existe et est unique.
\end{theorem}

\begin{proof}
  \emph{Unicité.} Si $g_1$ et $g_2$ satisfont la propriété, alors
  $\inner{\vec{u},g_1(\vec{v})-g_2(\vec{v})}=0$ pour tous
  $\vec{u},\vec{v}$. En prenant $\vec{u}=g_1(\vec{v})-g_2(\vec{v})$,
  on obtient $g_1(\vec{v})=g_2(\vec{v})$ pour tout $\vec{v}$.

  \emph{Existence.} Soit $(\vece{1},\ldots,\vece{n})$ une base
  orthonormée de $E$. Posons $A=\Mat_\cB(f)$. Définissons $g$ par
  $\Mat_\cB(g)=\hermit{A}$ (ou $\transpose{A}$ dans le cas réel). Alors~:
  \[
    \inner{f(\vec{u}),\vec{v}}=\hermit{(AX)}Y=\hermit{X}\hermit{A}Y
    =\inner{\vec{u},g(\vec{v})},
  \]
  où $X,Y$ sont les vecteurs de coordonnées dans la base orthonormée.
  Donc $g=f^*$.
\end{proof}

\begin{proposition}{Propriétés de l'adjoint}{prop:adjoint-prop}
  Soient $f,g\in\End(E)$ et $\lambda\in\K$. Alors~:
  \begin{enumerate}[label=(\roman*)]
    \item $(f+g)^*=f^*+g^*$.
    \item $(\lambda f)^*=\bar{\lambda}\,f^*$.
    \item $(f\circ g)^*=g^*\circ f^*$.
    \item $(f^*)^*=f$.
    \item $\Ker(f^*)=(\im f)^\perp$ et $\im(f^*)=(\Ker f)^\perp$.
  \end{enumerate}
\end{proposition}

\begin{proof}
  Les propriétés (i)--(iv) se déduisent de la définition par unicité de
  l'adjoint.

  Pour (v)~: $\vec{v}\in\Ker(f^*)$ si et seulement si $f^*(\vec{v})
  =\vzero$, si et seulement si $\inner{\vec{u},f^*(\vec{v})}=0$ pour
  tout $\vec{u}$, si et seulement si $\inner{f(\vec{u}),\vec{v}}=0$
  pour tout $\vec{u}$, si et seulement si $\vec{v}\perp\im(f)$, i.e.\
  $\vec{v}\in(\im f)^\perp$. La seconde égalité s'obtient en
  appliquant la première à $f^*$ et en utilisant $(f^*)^*=f$.
\end{proof}

\begin{remark}{Matrice de l'adjoint}{rem:mat-adjoint}
  \index{adjoint!matrice}
  Dans une base orthonormée $\cB$, si $A=\Mat_\cB(f)$, alors
  $\Mat_\cB(f^*)=\hermit{A}$ (ou $\transpose{A}$ si $\K=\R$). C'est
  pourquoi on écrit souvent $\hermit{A}$ pour la matrice adjointe.
\end{remark}

\begin{definition}{Endomorphismes spéciaux}{def:endomorphismes-speciaux}
  \index{endomorphisme auto-adjoint}\index{endomorphisme unitaire}
  \index{matrice symétrique}\index{matrice hermitienne}
  Soit $f\in\End(E)$ sur un espace préhilbertien de dimension finie.
  \begin{enumerate}[label=(\roman*)]
    \item $f$ est \emph{auto-adjoint} (ou \emph{symétrique} si $\K=\R$,
      \emph{hermitien} si $\K=\C$) si $f^*=f$. En base orthonormée,
      cela correspond à $A=\hermit{A}$.
    \item $f$ est \emph{orthogonal} (si $\K=\R$) ou \emph{unitaire}
      (si $\K=\C$) si $f^*\circ f=\Id$, i.e.\
      $f$ est une isométrie.
    \item $f$ est \emph{normal} si $f^*\circ f=f\circ f^*$.
  \end{enumerate}
\end{definition}

% ============================================================
\section{Applications}\label{sec:applications-ps}
% ============================================================

\subsection{Approximation aux moindres carrés}\label{subsec:moindres-carres-ch6}
\index{moindres carrés}

Considérons un système linéaire $A\vec{x}=\vec{b}$ avec
$A\in\Mat_{m\times n}(\R)$ et $m>n$ (système surdéterminé, plus
d'équations que d'inconnues). En général, ce système n'a pas de
solution. On cherche alors $\vec{x}^*\in\R^n$ qui minimise
$\norm{A\vec{x}-\vec{b}}^2$, c'est-à-dire le vecteur $\vec{x}^*$ tel
que $A\vec{x}^*$ soit la projection orthogonale de $\vec{b}$ sur
$\im(A)$.

\begin{proposition}{Équations normales}{prop:eq-normales}
  \index{équations normales}
  Le vecteur $\vec{x}^*$ minimisant $\norm{A\vec{x}-\vec{b}}^2$ est
  l'unique solution (si $A$ est de rang $n$) de~:
  \[
    \boxed{\transpose{A}\,A\,\vec{x}^*=\transpose{A}\,\vec{b}.}
  \]
\end{proposition}

\begin{proof}
  La condition d'optimalité est $\vec{b}-A\vec{x}^*\perp\im(A)$, i.e.\
  $\inner{A\vec{y},\vec{b}-A\vec{x}^*}=0$ pour tout $\vec{y}\in\R^n$.
  Cela donne $\transpose{\vec{y}}\transpose{A}(\vec{b}-A\vec{x}^*)=0$
  pour tout $\vec{y}$, d'où
  $\transpose{A}A\vec{x}^*=\transpose{A}\vec{b}$.
\end{proof}

\subsection{Aperçu de la factorisation QR}\label{subsec:QR}
\index{factorisation QR}

Le procédé de Gram-Schmidt fournit une factorisation matricielle
importante. Si $A\in\Mat_{m\times n}(\R)$ est de rang $n$ (colonnes
linéairement indépendantes), en appliquant Gram-Schmidt aux colonnes
$\vec{a}_1,\ldots,\vec{a}_n$ de $A$, on obtient une famille
orthonormée $\vec{q}_1,\ldots,\vec{q}_n$. Ceci conduit à~:
\[
  A=QR,
\]
où $Q\in\Mat_{m\times n}(\R)$ a des colonnes orthonormées
($\transpose{Q}Q=I_n$) et $R\in\Mat_n(\R)$ est triangulaire supérieure
avec des éléments diagonaux strictement positifs. La factorisation QR
est un outil fondamental en calcul numérique pour résoudre les systèmes
aux moindres carrés et calculer les valeurs propres.

\subsection{Aperçu des séries de Fourier}\label{subsec:fourier}
\index{séries de Fourier}

Sur l'espace $\mathcal{C}([-\pi,\pi],\R)$ muni du produit scalaire
$\inner{f,g}=\frac{1}{\pi}\int_{-\pi}^{\pi}f(t)g(t)\,dt$, la famille~:
\[
  \left(\frac{1}{\sqrt{2}},\;\cos(t),\;\sin(t),\;\cos(2t),\;
  \sin(2t),\;\ldots\right)
\]
est orthonormée. La projection orthogonale d'une fonction $f$ sur le
sous-espace engendré par les $2N+1$ premières fonctions est le
polynôme trigonométrique~:
\[
  S_N f(t)=\frac{a_0}{2}+\sum_{k=1}^N\bigl(a_k\cos(kt)
  +b_k\sin(kt)\bigr),
\]
où $a_k=\frac{1}{\pi}\int_{-\pi}^{\pi}f(t)\cos(kt)\,dt$ et
$b_k=\frac{1}{\pi}\int_{-\pi}^{\pi}f(t)\sin(kt)\,dt$ sont les
\emph{coefficients de Fourier} de $f$. L'inégalité de Bessel assure~:
\[
  \frac{a_0^2}{2}+\sum_{k=1}^N(a_k^2+b_k^2)
  \leq\frac{1}{\pi}\int_{-\pi}^{\pi}f(t)^2\,dt,
\]
et l'identité de Parseval (pour les séries de Fourier) affirme que
l'égalité a lieu à la limite $N\to+\infty$.

% ============================================================
\section{Illustrations}\label{sec:illustrations-ps}
% ============================================================

\begin{figure}[ht]
\centering
\begin{tikzpicture}[scale=1.8, >=Stealth]
  % Axes
  \draw[gray!30, very thin] (-0.5,-0.5) grid (4.5,3.5);
  \draw[->] (-0.5,0) -- (4.7,0) node[right] {$x$};
  \draw[->] (0,-0.5) -- (0,3.7) node[above] {$y$};

  % Subspace F (line through origin)
  \draw[thick, blue!50, dashed] (-0.3,-0.15) -- (4.5,2.25);
  \node[blue!70!black] at (4.2,2.5) {$F$};

  % Vector v
  \coordinate (v) at (2,3);
  \draw[->, thick, red!70!black] (0,0) -- (v) node[above right] {$\vec{v}$};

  % Projection
  \coordinate (p) at (2.769,1.385);
  \draw[->, thick, blue!70!black] (0,0) -- (p)
    node[below right] {$\proj_F(\vec{v})$};

  % Orthogonal component
  \draw[->, thick, green!50!black, dashed] (p) -- (v)
    node[midway, right=3pt] {$\vec{v}-\proj_F(\vec{v})$};

  % Right angle
  \draw[thick] ($(p)+(0.15,0.075)$) -- +(-0.075,0.15) -- +(-0.225,0.075);

  % Distance annotation
  \draw[<->, thin, purple] (2.1,3.05) -- node[above, sloped, font=\small]
    {$d(\vec{v},F)$} (2.869,1.44);
\end{tikzpicture}
\caption{Projection orthogonale de $\vec{v}$ sur un sous-espace $F$~:
  le projeté $\proj_F(\vec{v})$ est le point de $F$ le plus proche de
  $\vec{v}$, et le vecteur $\vec{v}-\proj_F(\vec{v})$ est orthogonal
  à~$F$.}
\label{fig:proj-ortho}
\end{figure}

\begin{figure}[ht]
\centering
\begin{tikzpicture}[scale=1.6, >=Stealth]
  % Original vectors
  \coordinate (v1) at (3,1);
  \coordinate (v2) at (1,2.5);

  % Gram-Schmidt result
  \coordinate (e1) at ({3/sqrt(10)},{1/sqrt(10)});  % v1 normalized
  % e2 direction: v2 - proj_{e1}(v2) e1
  % <e1,v2> = (3*1+1*2.5)/sqrt(10) = 5.5/sqrt(10)
  % proj = 5.5/10 * (3,1) = (1.65, 0.55)
  % w2 = (1-1.65, 2.5-0.55) = (-0.65, 1.95)
  \coordinate (w2) at (-0.65, 1.95);

  % Draw original vectors
  \draw[->, thick, red!60!black] (0,0) -- (v1)
    node[below right] {$\vec{v}_1$};
  \draw[->, thick, red!60!black] (0,0) -- (v2)
    node[above left] {$\vec{v}_2$};

  % Draw orthonormal vectors (scaled for visibility)
  \draw[->, very thick, blue!70!black] (0,0) -- ($2.5*(e1)$)
    node[below] {$\vece{1}$};

  \pgfmathsetmacro{\wnorm}{sqrt(0.65*0.65+1.95*1.95)}
  \draw[->, very thick, blue!70!black] (0,0) --
    ({-0.65/\wnorm*2.5},{1.95/\wnorm*2.5})
    node[above left] {$\vece{2}$};

  % Show projection of v2 onto e1
  \draw[->, thin, gray, dashed] (0,0) -- (1.65,0.55)
    node[below, font=\small] {$\proj_{\vece{1}}(\vec{v}_2)$};
  \draw[thin, gray, dashed] (v2) -- (1.65,0.55);

  % w2 arrow
  \draw[->, thick, green!50!black] (1.65,0.55) -- (v2)
    node[midway, right=2pt, font=\small] {$\vec{w}_2$};

  % Right angle mark
  \draw[thick] ($(1.65,0.55)+(0.12,0.04)$) -- +(0.04,-0.12) -- +(-0.08,-0.08);
\end{tikzpicture}
\caption{Procédé de Gram-Schmidt en dimension 2~: à partir de
  $\vec{v}_1$ et $\vec{v}_2$, on construit $\vece{1}$ (normalisation
  de $\vec{v}_1$) puis $\vec{w}_2=\vec{v}_2-\proj_{\vece{1}}(\vec{v}_2)$,
  et enfin $\vece{2}$ (normalisation de $\vec{w}_2$).}
\label{fig:gram-schmidt}
\end{figure}

\begin{figure}[ht]
\centering
\begin{tikzpicture}[scale=1.8, >=Stealth]
  % Vectors
  \coordinate (O) at (0,0);
  \coordinate (u) at (3,0.5);
  \coordinate (v) at (1.5,2.5);

  \draw[->, thick, blue!70!black] (O) -- (u) node[right] {$\vec{u}$};
  \draw[->, thick, red!70!black] (O) -- (v) node[above] {$\vec{v}$};

  % Angle arc
  \pic[draw, thick, "$\theta$" font=\normalsize, angle eccentricity=1.4,
    angle radius=1.2cm] {angle=u--O--v};

  % Formula
  \node[anchor=west, text width=5cm] at (3.5,1.5) {%
    $\cos\theta=\dfrac{\inner{\vec{u},\vec{v}}}%
    {\norm{\vec{u}}\cdot\norm{\vec{v}}}$};
\end{tikzpicture}
\caption{Angle $\theta$ entre deux vecteurs~: le cosinus de l'angle est
  donné par le rapport du produit scalaire au produit des normes.}
\label{fig:angle-vecteurs}
\end{figure}

% ============================================================
\section{Exercices}\label{sec:exos-ps}
% ============================================================

\begin{exercise}{\basic{} Produit scalaire et norme sur $\R^2$}{exo:ps-R2}
  Soit $E=\R^2$ muni du produit scalaire canonique.
  \begin{enumerate}[label=(\alph*)]
    \item Calculer $\inner{\vec{u},\vec{v}}$, $\norm{\vec{u}}$,
      $\norm{\vec{v}}$ et l'angle $\theta$ entre $\vec{u}$ et $\vec{v}$,
      pour $\vec{u}=(1,2)$ et $\vec{v}=(3,-1)$.
    \item Trouver tous les vecteurs de $\R^2$ de norme $1$ orthogonaux
      à $\vec{u}=(3,4)$.
  \end{enumerate}
\end{exercise}

\begin{exercise}{\basic{} Vérification des axiomes}{exo:verif-ps}
  Déterminer si les applications suivantes définissent un produit
  scalaire sur $\R^2$~:
  \begin{enumerate}[label=(\alph*)]
    \item $\vphi(\vec{x},\vec{y})=x_1 y_1-x_1 y_2-x_2 y_1+4x_2 y_2$.
    \item $\psi(\vec{x},\vec{y})=x_1 y_1-x_1 y_2-x_2 y_1+x_2 y_2$.
    \item $\chi(\vec{x},\vec{y})=2x_1 y_1+x_1 y_2+x_2 y_1+2x_2 y_2$.
  \end{enumerate}
  \emph{Indication~:} vérifier la bilinéarité, la symétrie et le
  caractère défini positif. Pour ce dernier, examiner la matrice
  associée.
\end{exercise}

\begin{exercise}{\basic{} Orthogonal d'un sous-espace}{exo:orthogonal}
  Dans $\R^3$ muni du produit scalaire canonique, soit
  $F=\Vect\left(\begin{pmatrix}1\\1\\0\end{pmatrix},
  \begin{pmatrix}0\\1\\1\end{pmatrix}\right)$.
  \begin{enumerate}[label=(\alph*)]
    \item Déterminer $F^\perp$.
    \item Vérifier que $\R^3=F\oplus F^\perp$.
    \item Calculer la projection orthogonale de
      $\vec{v}=\begin{pmatrix}1\\2\\3\end{pmatrix}$ sur $F$.
  \end{enumerate}
\end{exercise}

\begin{exercise}{\basic{} Gram-Schmidt dans $\R^2$}{exo:gs-R2}
  Appliquer le procédé de Gram-Schmidt à la base
  $\vec{v}_1=(3,4)$, $\vec{v}_2=(1,0)$ de $\R^2$ muni du produit
  scalaire canonique.
\end{exercise}

\begin{exercise}{\basic{} Inégalité de Cauchy-Schwarz}{exo:cauchy-schwarz-app}
  \begin{enumerate}[label=(\alph*)]
    \item En appliquant l'inégalité de Cauchy-Schwarz dans $\R^n$ avec
      des vecteurs bien choisis, montrer que pour tous réels
      $a_1,\ldots,a_n$~:
      \[
        \left(\sum_{i=1}^n a_i\right)^2\leq n\sum_{i=1}^n a_i^2.
      \]
    \item En appliquant Cauchy-Schwarz sur
      $\mathcal{C}([0,1],\R)$, montrer que pour toute
      $f\in\mathcal{C}([0,1],\R)$~:
      \[
        \left(\int_0^1 f(t)\,dt\right)^2
        \leq\int_0^1 f(t)^2\,dt.
      \]
  \end{enumerate}
\end{exercise}

\begin{exercise}{\intermediate{} Matrice de Gram et produit scalaire}{exo:gram}
  \index{matrice de Gram}
  Soit $(\vec{v}_1,\ldots,\vec{v}_p)$ une famille de vecteurs d'un
  espace préhilbertien $E$. La \emph{matrice de Gram} est la matrice
  $G\in\Mat_p(\K)$ définie par $G_{ij}=\inner{\vec{v}_i,\vec{v}_j}$.
  \begin{enumerate}[label=(\alph*)]
    \item Montrer que $G$ est hermitienne (symétrique si $\K=\R$) et
      semi-définie positive.
    \item Montrer que $\det(G)=0$ si et seulement si la famille
      $(\vec{v}_i)$ est liée.
    \item Calculer la matrice de Gram de la famille
      $\vec{v}_1=(1,0,1)$, $\vec{v}_2=(1,1,0)$ dans $\R^3$.
  \end{enumerate}
\end{exercise}

\begin{exercise}{\intermediate{} Projection orthogonale et distance}{exo:proj-dist}
  Dans $\R^4$ muni du produit scalaire canonique, soit
  $F=\Vect\left(\begin{pmatrix}1\\0\\1\\0\end{pmatrix},
  \begin{pmatrix}0\\1\\0\\1\end{pmatrix}\right)$.
  \begin{enumerate}[label=(\alph*)]
    \item Construire une base orthonormée de $F$.
    \item Calculer la projection orthogonale de
      $\vec{v}=(1,2,3,4)$ sur $F$.
    \item En déduire $d(\vec{v},F)$.
  \end{enumerate}
\end{exercise}

\begin{exercise}{\intermediate{} Gram-Schmidt et factorisation QR}{exo:QR}
  Soit $A=\begin{pmatrix}1&1\\1&0\\0&1\end{pmatrix}$.
  \begin{enumerate}[label=(\alph*)]
    \item Appliquer le procédé de Gram-Schmidt aux colonnes de $A$ pour
      obtenir une famille orthonormée $(\vec{q}_1,\vec{q}_2)$.
    \item Écrire la factorisation $A=QR$ où $Q$ a des colonnes
      orthonormées et $R$ est triangulaire supérieure.
    \item Utiliser cette factorisation pour résoudre le système aux
      moindres carrés $A\vec{x}\approx\vec{b}$ avec
      $\vec{b}=(1,0,0)$.
  \end{enumerate}
\end{exercise}

\begin{exercise}{\intermediate{} Matrices orthogonales}{exo:mat-ortho}
  \begin{enumerate}[label=(\alph*)]
    \item Montrer que le produit de deux matrices orthogonales est
      orthogonal.
    \item Montrer que l'inverse d'une matrice orthogonale est orthogonale.
    \item Déterminer toutes les matrices de $\Orth(2)$ et les classer
      en rotations et réflexions.
    \item La matrice
      $\frac{1}{3}\begin{pmatrix}2&-1&2\\2&2&-1\\-1&2&2\end{pmatrix}$
      est-elle orthogonale~? Si oui, est-elle dans $\SO(3)$~?
  \end{enumerate}
\end{exercise}

\begin{exercise}{\intermediate{} Endomorphisme auto-adjoint}{exo:auto-adjoint}
  Soit $E=\R^3$ muni du produit scalaire canonique et
  $f\colon E\to E$ défini par $f(x,y,z)=(2x+y,\,x+2y,\,z)$.
  \begin{enumerate}[label=(\alph*)]
    \item Écrire la matrice $A$ de $f$ dans la base canonique.
    \item Vérifier que $f$ est auto-adjoint.
    \item Diagonaliser $f$ en base orthonormée.
  \end{enumerate}
\end{exercise}

\begin{exercise}{\challenging{} Polynômes orthogonaux de Legendre}{exo:legendre}
  \index{polynômes de Legendre}
  Sur $E=\R_2[X]$ (polynômes de degré $\leq 2$), on définit le produit
  scalaire~:
  \[
    \inner{P,Q}=\int_{-1}^1 P(t)\,Q(t)\,dt.
  \]
  \begin{enumerate}[label=(\alph*)]
    \item Vérifier que c'est bien un produit scalaire.
    \item Appliquer le procédé de Gram-Schmidt à la base $(1,X,X^2)$
      pour obtenir une base orthogonale $(P_0,P_1,P_2)$.
    \item Normaliser pour obtenir une base orthonormée.
    \item Les polynômes $P_0,P_1,P_2$ obtenus sont (à un facteur
      multiplicatif près) les premiers polynômes de Legendre. Vérifier
      que $P_2(t)$ est proportionnel à $\frac{1}{2}(3t^2-1)$.
  \end{enumerate}
\end{exercise}

\begin{exercise}{\challenging{} Moindres carrés~: régression linéaire}{exo:moindres-carres-ch6}
  \index{régression linéaire}
  On dispose des données $(x_i,y_i)_{1\leq i\leq 4}$ suivantes~:
  $(1,2)$, $(2,3)$, $(3,5)$, $(4,4)$.
  On cherche la droite $y=a+bx$ qui minimise
  $\sum_{i=1}^4(y_i-a-bx_i)^2$.
  \begin{enumerate}[label=(\alph*)]
    \item Formuler ce problème comme un système surdéterminé
      $A\vec{c}=\vec{y}$ avec $\vec{c}=(a,b)$.
    \item Écrire et résoudre les équations normales.
    \item Calculer l'erreur résiduelle
      $\norm{A\vec{c}^*-\vec{y}}$.
  \end{enumerate}
\end{exercise}

\begin{exercise}{\challenging{} Isométries du plan}{exo:isometries}
  \index{isométrie}
  Soit $f\colon\R^2\to\R^2$ un endomorphisme dont la matrice dans la base
  canonique est $A\in\Mat_2(\R)$.
  \begin{enumerate}[label=(\alph*)]
    \item Montrer que $f$ est une isométrie si et seulement si
      $A\in\Orth(2)$.
    \item En déduire que toute isométrie de $\R^2$ est soit une rotation,
      soit une réflexion.
    \item Soit $f$ la réflexion par rapport à la droite passant par
      l'origine et faisant un angle $\alpha$ avec l'axe des abscisses.
      Montrer que~:
      \[
        \Mat(f)=\begin{pmatrix}\cos 2\alpha&\sin 2\alpha\\
        \sin 2\alpha&-\cos 2\alpha\end{pmatrix}.
      \]
    \item Montrer que la composée de deux réflexions par rapport à des
      droites passant par l'origine est une rotation, et déterminer son
      angle en fonction des angles des deux droites.
  \end{enumerate}
\end{exercise}

\begin{exercise}{\challenging{} Adjoint et noyau}{exo:adjoint-noyau}
  Soit $f\in\End(E)$ sur un espace euclidien de dimension $n$.
  \begin{enumerate}[label=(\alph*)]
    \item Montrer que $\rg(f)=\rg(f^*)=\rg(\transpose{A}A)$ où
      $A=\Mat(f)$ dans une base orthonormée.
    \item Montrer que $\Ker(f^*\circ f)=\Ker(f)$.

      \emph{Indication~:} considérer
      $\inner{f^*(f(\vec{v})),\vec{v}}=\norm{f(\vec{v})}^2$.
    \item En déduire que $\rg(\transpose{A}A)=\rg(A)$ pour toute
      matrice réelle $A$.
  \end{enumerate}
\end{exercise}

\begin{exercise}{\challenging{} Approximation de Fourier}{exo:fourier-approx}
  \index{approximation de Fourier}
  Sur $\mathcal{C}([-\pi,\pi],\R)$ muni du produit scalaire
  $\inner{f,g}=\frac{1}{\pi}\int_{-\pi}^{\pi}f(t)g(t)\,dt$, on
  considère $f(t)=t$.
  \begin{enumerate}[label=(\alph*)]
    \item Calculer $a_0$, $a_k=\frac{1}{\pi}\int_{-\pi}^{\pi}
      t\cos(kt)\,dt$ et $b_k=\frac{1}{\pi}\int_{-\pi}^{\pi}
      t\sin(kt)\,dt$ pour $k\geq 1$.
    \item Écrire le polynôme trigonométrique $S_N f(t)$ de meilleure
      approximation de $f$ à l'ordre $N$.
    \item En utilisant l'identité de Parseval (à la limite),
      montrer que $\sum_{k=1}^{+\infty}\frac{1}{k^2}
      =\frac{\pi^2}{6}$.
  \end{enumerate}
\end{exercise}

% ============================================================
\section*{Résumé du chapitre}
\addcontentsline{toc}{section}{Résumé du chapitre}
% ============================================================

\begin{framed}
\textbf{Ce qu'il faut retenir.}

\begin{enumerate}[label=\textbf{\arabic*.},leftmargin=2em]
  \item \textbf{Forme bilinéaire et forme quadratique.} Une forme
    bilinéaire symétrique $\vphi$ est représentée par une matrice
    symétrique. La forme quadratique associée est
    $q(\vec{v})=\vphi(\vec{v},\vec{v})$, et on peut retrouver $\vphi$
    à partir de $q$ par polarisation. La loi d'inertie de Sylvester
    assure que la signature $(p,q)$ est un invariant.

  \item \textbf{Produit scalaire.} Un produit scalaire est une forme
    bilinéaire symétrique (cas réel) ou sesquilinéaire hermitienne
    (cas complexe), définie positive. Il permet de définir norme,
    distance et angle.

  \item \textbf{Cauchy-Schwarz.}
    $\abs{\inner{\vec{u},\vec{v}}}
    \leq\norm{\vec{u}}\cdot\norm{\vec{v}}$, avec égalité ssi
    $\vec{u}$ et $\vec{v}$ sont colinéaires. Conséquence~: inégalité
    triangulaire.

  \item \textbf{Orthogonalité.} $\vec{u}\perp\vec{v}$ ssi
    $\inner{\vec{u},\vec{v}}=0$. Toute famille orthogonale de vecteurs
    non nuls est libre. En dimension finie, $E=F\oplus F^\perp$ et
    $\dim F^\perp=\dim E-\dim F$.

  \item \textbf{Projection orthogonale.} $\proj_F(\vec{v})$ est l'unique
    élément de $F$ le plus proche de $\vec{v}$. En base orthonormée
    de $F$~:
    $\proj_F(\vec{v})=\sum_i\inner{\vece{i},\vec{v}}\,\vece{i}$.

  \item \textbf{Gram-Schmidt.} Algorithme d'orthonormalisation~: à
    partir d'une famille libre, on construit une famille orthonormée
    engendrant les mêmes sous-espaces. Tout espace préhilbertien de
    dimension finie admet une base orthonormée.

  \item \textbf{Parseval et Bessel.} En base orthonormée complète~:
    $\norm{\vec{v}}^2=\sum_i\abs{\inner{\vece{i},\vec{v}}}^2$
    (Parseval). Pour une famille orthonormée partielle, on a
    l'inégalité de Bessel $\leq$.

  \item \textbf{Matrices orthogonales et unitaires.}
    $Q\in\Orth(n)\Leftrightarrow\transpose{Q}Q=I_n$~: les colonnes
    forment une base orthonormée. $\det(Q)=\pm 1$~;
    $\SO(n)=\setcond{Q\in\Orth(n)}{\det Q=1}$.

  \item \textbf{Adjoint.} $f^*$ est l'unique endomorphisme tel que
    $\inner{f(\vec{u}),\vec{v}}=\inner{\vec{u},f^*(\vec{v})}$. En
    base orthonormée, $\Mat(f^*)=\hermit{\Mat(f)}$.
    $\Ker(f^*)=(\im f)^\perp$.

  \item \textbf{Applications.}
    \begin{itemize}
      \item Moindres carrés~: $\transpose{A}A\vec{x}^*
        =\transpose{A}\vec{b}$.
      \item Factorisation QR~: $A=QR$ (Gram-Schmidt sur les colonnes).
      \item Séries de Fourier~: projection orthogonale sur les
        sous-espaces trigonométriques.
    \end{itemize}
\end{enumerate}
\end{framed}
